\documentclass[12pt,a4paper]{article}
\usepackage[a4paper,margin=1.8cm]{geometry}
\usepackage[T2A]{fontenc}
\usepackage[utf8]{inputenc}
\usepackage[russian]{babel}
\usepackage{amsmath,amssymb,mathtools}
\usepackage{array,booktabs,longtable}
\usepackage{xcolor}
\usepackage{hyperref}
\usepackage{tikz}
\hypersetup{colorlinks=true,linkcolor=blue!55!black,urlcolor=blue!55!black}
\setlength{\parindent}{0pt}
\setlength{\parskip}{5pt}
\renewcommand{\arraystretch}{1.18}
\newcommand{\err}{\textcolor{red}{Ошибка:}}
\newcommand{\idea}{\textcolor{blue!60!black}{Ключевая идея:}}
\newcommand{\result}{\textcolor{teal!60!black}{Итог:}}
\newcommand{\tasktitle}[1]{\phantomsection\label{task:#1}\section*{Задание №#1}}

\begin{document}

\hypersetup{pageanchor=false}
\begin{titlepage}
\centering
\vspace*{0.22\textheight}
{\LARGE\bfseries Ревью домашней работы\par}
\vspace{1.6cm}
{\large Автор: Лёвин Артём Александрович\par}
\vspace{0.5cm}
{\large 3 октября 2026 г.\par}
\vfill
\begin{tikzpicture}[scale=1]
\draw[blue!45!black, line width=0.7pt] (0,0) .. controls (2,0.65) and (4,-0.65) .. (6,0);
\end{tikzpicture}
\end{titlepage}
\hypersetup{pageanchor=true}

\newpage
\section*{Сводная таблица}
\small
\begin{longtable}{>{\centering\arraybackslash}p{0.8cm} p{3.0cm} p{5.0cm} p{2.6cm} p{2.9cm}}
\toprule
№ & Тема & Суть ошибки & Ваш ответ & Правильный ответ \\
\midrule
\endfirsthead
\toprule
№ & Тема & Суть ошибки & Ваш ответ & Правильный ответ \\
\midrule
\endhead
\hyperref[task:4]{4} & Уравнение с модулем и дробями & В ответ включено значение, при котором знаменатель равен нулю; дополнительно неверно решено условие равенства дроби нулю. & $8;\,-1;\,-\frac23$ & $8;\,-\frac23$ \\
\addlinespace
\hyperref[task:5]{5} & Сумма модулей & На промежутке $x<-1$ неверно сложены постоянные слагаемые, поэтому получен неверный корень. & По записи: $-2{,}5$ и $4$ & $-3;\,4$ \\
\addlinespace
\hyperref[task:7]{7} & Сумма модулей, тождество на промежутке & Равенство $0=0$ ошибочно принято за отсутствие корней, хотя оно означает, что подходит весь рассматриваемый промежуток. & $-2$ & $[-2;1]$ \\
\addlinespace
\hyperref[task:8]{8} & Разность модулей, тождество на промежутке & Тождество на промежутке $x<-2$ ошибочно отброшено как «нет корней». & $-2$ & $(-\infty;-2]$ \\
\addlinespace
\hyperref[task:10]{10} & Вложенный модуль & Решение не доведено до полного набора корней; в первой ветви потерян уже найденный корень $-2$. & Записан $6$, вторая ветвь не завершена & $-2;\,0;\,4;\,6$ \\
\bottomrule
\end{longtable}
\normalsize

\newpage
\vspace*{0.32\textheight}
\begin{center}
{\LARGE\bfseries Разбор ошибок}
\end{center}
\vfill
\newpage

\tasktitle{4}
\textbf{Текст задания.} Решите уравнение
\[
\left|\frac{2x-3}{x+1}\right|=\frac{x+5}{x+1}.
\]

\textbf{Ваше решение.} В записи рассмотрены два случая для выражения под модулем. В первом случае получено
\[
\frac{2x-3}{x+1}=\frac{x+5}{x+1},
\]
откуда после переноса получено
\[
\frac{x-8}{x+1}=0.
\]
Далее найдено $x=8$, но дополнительно из знаменателя записано $x=-1$ и это значение включено в ответ. Во втором случае получено значение $x=-\frac23$.

\textbf{Ваш ответ.} $8;\,-1;\,-\frac23$.

\err{} в ответ включено $x=-1$, хотя при этом значении исходное выражение не определено. Кроме того, дробь равна нулю только тогда, когда её числитель равен нулю, а знаменатель при этом не равен нулю.

\textbf{Где именно возникает ошибка.} Последний корректный шаг в первой ветви:
\[
\frac{x-8}{x+1}=0.
\]
Первый неверный шаг --- попытка получить отдельный корень из условия $x+1=0$. Это условие не даёт корень уравнения, а, наоборот, запрещает значение $x=-1$.

\textbf{Почему это неверно.} Для дроби
\[
\frac{A}{B}
\]
равенство нулю возможно только при двух одновременных условиях:
\[
A=0,\qquad B\ne0.
\]
Если знаменатель равен нулю, дробь не существует. Поэтому из
\[
\frac{x-8}{x+1}=0
\]
следует $x=8$ и обязательно сохраняется ограничение $x\ne-1$. Значение $x=-1$ не может быть корнем ни этой ветви, ни исходного уравнения.

\idea{} перед преобразованиями дробного уравнения нужно выписать область допустимых значений, а при равенстве дроби нулю обнулять только числитель.

\textbf{Исправление от последнего правильного шага.} Из
\[
\frac{x-8}{x+1}=0
\]
получаем
\[
x-8=0,\qquad x+1\ne0.
\]
Значит,
\[
x=8.
\]
Вторая ветвь даёт
\[
\frac{2x-3}{x+1}=-\frac{x+5}{x+1},
\]
то есть
\[
2x-3=-x-5,
\]
\[
3x=-2,
\]
\[
x=-\frac23.
\]
Оба значения не равны $-1$.

\textbf{Полное правильное решение.} Сначала фиксируем область допустимых значений:
\[
x\ne-1.
\]
Правая часть исходного равенства должна быть неотрицательной для любого корня. Рассмотрим два стандартных случая.

Первый случай:
\[
\frac{2x-3}{x+1}=\frac{x+5}{x+1}.
\]
Так как $x\ne-1$, умножаем обе части на $x+1$:
\[
2x-3=x+5,
\]
\[
x=8.
\]
Проверка правой части:
\[
\frac{8+5}{8+1}=\frac{13}{9}>0.
\]
Значение подходит.

Второй случай:
\[
\frac{2x-3}{x+1}=-\frac{x+5}{x+1}.
\]
При $x\ne-1$:
\[
2x-3=-x-5,
\]
\[
3x=-2,
\]
\[
x=-\frac23.
\]
Проверка правой части:
\[
\frac{-\frac23+5}{-\frac23+1}=13>0.
\]
Значение подходит.

\textbf{Проверка.} Подстановка $x=8$ и $x=-\frac23$ в исходное уравнение даёт верные равенства. Значение $x=-1$ подставлять нельзя, потому что знаменатель обращается в ноль.

\result{} $x=8$ или $x=-\frac23$.

\subsection*{Задача на закрепление}
Решите уравнение
\[
\left|\frac{3x-2}{x+2}\right|=\frac{x+6}{x+2}.
\]

\newpage
\tasktitle{5}
\textbf{Текст задания.} Решите уравнение
\[
|x-2|+|x+1|=7.
\]

\textbf{Ваше решение.} Правильно найдены критические точки $x=-1$ и $x=2$ и произведено разбиение числовой прямой на три промежутка. Для среднего промежутка сделан верный вывод об отсутствии решений, а для правого промежутка получено $x=4$. На левом промежутке записано
\[
-(x-2)-(x+1)-7=0,
\]
но далее из этой строки получено $-2x=5$ и затем $x=-2{,}5$.

\textbf{Ваш ответ.} По записи получены значения $-2{,}5$ и $4$; отдельная итоговая строка ответа не оформлена.

\err{} при приведении подобных слагаемых на промежутке $x<-1$ неверно сложены постоянные числа.

\textbf{Где именно возникает ошибка.} Последний корректный шаг:
\[
-(x-2)-(x+1)-7=0.
\]
После раскрытия скобок должно получиться
\[
-x+2-x-1-7=0.
\]
Первый неверный переход --- замена этой строки на равенство, эквивалентное $-2x=5$.

\textbf{Почему это неверно.} Постоянные слагаемые складываются так:
\[
2-1-7=-6.
\]
Поэтому левая часть равна
\[
-2x-6,
\]
а не выражению, которое приводит к $-2x=5$. Из-за одной арифметической ошибки корень левой ветви сдвинулся с $-3$ на $-2{,}5$.

\idea{} после раскрытия модулей полезно отдельно собрать коэффициенты при $x$ и отдельно постоянные слагаемые, а затем проверить найденный корень принадлежностью выбранному промежутку.

\textbf{Исправление от последнего правильного шага.} На промежутке $x<-1$:
\[
-(x-2)-(x+1)=7,
\]
\[
-x+2-x-1=7,
\]
\[
-2x+1=7,
\]
\[
-2x=6,
\]
\[
x=-3.
\]
Условие $x<-1$ выполнено.

\textbf{Полное правильное решение.} Критические точки:
\[
x=-1,\qquad x=2.
\]

Если $x<-1$, оба выражения под модулями отрицательны:
\[
-(x-2)-(x+1)=7,
\]
\[
-2x+1=7,
\]
\[
x=-3.
\]
Значение принадлежит промежутку.

Если $-1\le x<2$, то
\[
-(x-2)+(x+1)=7,
\]
\[
3=7.
\]
Решений на этом промежутке нет.

Если $x\ge2$, то
\[
(x-2)+(x+1)=7,
\]
\[
2x-1=7,
\]
\[
2x=8,
\]
\[
x=4.
\]
Значение принадлежит промежутку.

\textbf{Проверка.} При $x=-3$:
\[
|-5|+|-2|=5+2=7.
\]
При $x=4$:
\[
|2|+|5|=2+5=7.
\]

\result{} $x=-3$ или $x=4$.

\subsection*{Задача на закрепление}
Решите уравнение
\[
|x-3|+|x+2|=9.
\]

\newpage
\tasktitle{7}
\textbf{Текст задания.} Решите уравнение
\[
|x-1|+|x+2|=3.
\]

\textbf{Ваше решение.} Правильно найдены критические точки $x=-2$ и $x=1$. На среднем промежутке после раскрытия модулей получено тождество $0=0$, но рядом сделан вывод «нет корней». В итоге в ответ оставлено только значение $-2$.

\textbf{Ваш ответ.} $-2$.

\err{} тождество $0=0$ истолковано как отсутствие решений. На самом деле тождество означает, что равенство выполняется при каждом значении $x$ из рассматриваемого промежутка.

\textbf{Где именно возникает ошибка.} Последний корректный шаг на промежутке $-2<x<1$:
\[
-(x-1)+(x+2)-3=0.
\]
После упрощения получается
\[
0=0.
\]
Первый неверный вывод --- «нет корней».

\textbf{Почему это неверно.} Равенство $0=0$ истинно всегда. При решении по промежуткам переменная уже ограничена условием выбранного промежутка. Значит, если после раскрытия модулей получается тождество, подходят все значения из этого промежутка. В данном случае фактически сумма расстояний от точки $x$ до точек $-2$ и $1$ равна расстоянию между самими точками, когда $x$ лежит между ними.

\idea{} нужно различать два принципиально разных результата: противоречие вроде $0=5$ означает отсутствие решений, а тождество $0=0$ означает, что подходят все значения текущего промежутка.

\textbf{Исправление от последнего правильного шага.} На промежутке $-2\le x\le1$:
\[
|x-1|+|x+2|=(1-x)+(x+2)=3.
\]
Равенство выполняется для каждого $x$ из этого промежутка. Поэтому весь отрезок $[-2;1]$ входит в множество решений.

\textbf{Полное правильное решение.} Рассмотрим три промежутка.

Если $x<-2$, то
\[
|x-1|+|x+2|=(1-x)+(-x-2)=-2x-1.
\]
Уравнение
\[
-2x-1=3
\]
даёт $x=-2$, но это значение не принадлежит открытому промежутку $x<-2$.

Если $-2\le x\le1$, то
\[
|x-1|+|x+2|=(1-x)+(x+2)=3.
\]
Получаем тождество, поэтому подходят все значения из $[-2;1]$.

Если $x>1$, то
\[
|x-1|+|x+2|=(x-1)+(x+2)=2x+1.
\]
Уравнение
\[
2x+1=3
\]
даёт $x=1$, но это значение не принадлежит открытому промежутку $x>1$. Оно уже входит в средний отрезок.

\textbf{Проверка.} Внутри отрезка, например при $x=0$:
\[
|-1|+|2|=1+2=3.
\]
Для любой точки между $-2$ и $1$ сумма расстояний до концов отрезка равна длине самого отрезка, то есть $3$.

\result{} все $x$ из отрезка $[-2;1]$.

\subsection*{Задача на закрепление}
Решите уравнение
\[
|x-4|+|x+1|=5.
\]

\newpage
\tasktitle{8}
\textbf{Текст задания.} Решите уравнение
\[
|x-4|-|x+2|=6.
\]

\textbf{Ваше решение.} Правильно найдены критические точки $x=-2$ и $x=4$. На промежутке $x<-2$ после раскрытия модулей получено тождество, но сделан вывод «нет корней». На среднем промежутке найдено $x=-2$, и именно это значение записано как итоговый ответ.

\textbf{Ваш ответ.} $-2$.

\err{} отброшен целый промежуток решений. Тождество на промежутке $x<-2$ означает, что каждое значение из этого промежутка удовлетворяет исходному уравнению.

\textbf{Где именно возникает ошибка.} На промежутке $x<-2$:
\[
-(x-4)+ (x+2)-6=0.
\]
После упрощения получается тождество. Первый неверный шаг --- вывод об отсутствии корней.

\textbf{Почему это неверно.} Для $x<-2$ имеем
\[
|x-4|=4-x,
\qquad
|x+2|=-x-2.
\]
Тогда
\[
|x-4|-|x+2|=(4-x)-(-x-2)=6.
\]
Переменная полностью сокращается, потому что разность расстояний до точек $4$ и $-2$ для любой точки слева от обеих точек постоянна и равна расстоянию между ними. Значит, подходят все $x<-2$. Граничная точка $x=-2$ также подходит.

\idea{} если после раскрытия модулей получается истинное равенство без переменной, нужно оставить весь текущий промежуток, а не искать отдельный корень.

\textbf{Исправление от последнего правильного шага.} На промежутке $x<-2$:
\[
(4-x)-(-x-2)=6,
\]
\[
6=6.
\]
Следовательно, подходят все $x<-2$. Отдельно $x=-2$ даёт
\[
|-6|-|0|=6,
\]
поэтому граница также входит в ответ.

\textbf{Полное правильное решение.} Если $x<-2$, то
\[
|x-4|-|x+2|=(4-x)-(-x-2)=6.
\]
Значит, весь промежуток $(-\infty;-2)$ состоит из решений.

Если $-2\le x<4$, то
\[
|x-4|-|x+2|=(4-x)-(x+2)=2-2x.
\]
Решаем:
\[
2-2x=6,
\]
\[
-2x=4,
\]
\[
x=-2.
\]
Граничное значение подходит.

Если $x\ge4$, то
\[
|x-4|-|x+2|=(x-4)-(x+2)=-6,
\]
а равенство $-6=6$ невозможно.

\textbf{Проверка.} Например, при $x=-5$:
\[
|-9|-|-3|=9-3=6.
\]
При $x=-2$:
\[
|-6|-|0|=6.
\]
При любом $x>-2$ вне граничной точки требуемое равенство либо даёт другой вид, либо не выполняется.

\result{} $x\le-2$, то есть $(-\infty;-2]$.

\subsection*{Задача на закрепление}
Решите уравнение
\[
|x-5|-|x+1|=6.
\]

\newpage
\tasktitle{10}
\textbf{Текст задания.} Решите уравнение
\[
\bigl||x-2|-3\bigr|=1.
\]

\textbf{Ваше решение.} Правильно выполнен первый шаг раскрытия внешнего модуля:
\[
|x-2|-3=1
\quad\text{или}\quad
|x-2|-3=-1.
\]
Отсюда правильно получены два уравнения
\[
|x-2|=4
\quad\text{и}\quad
|x-2|=2.
\]
Для первого уравнения в записи появляются значения $x=6$ и $x=-2$, однако далее сохраняется только $x=6$. Второе уравнение начато, но решение на представленной странице не доведено до окончательных корней и итогового ответа.

\textbf{Ваш ответ.} Полный ответ не указан; в завершённой части записи оставлено значение $6$.

\err{} один уже найденный корень $-2$ потерян, а вторая ветвь решения не завершена, поэтому полный набор корней не получен.

\textbf{Где именно возникает ошибка.} До системы
\[
|x-2|=4
\quad\text{или}\quad
|x-2|=2
\]
решение полностью корректно. В первой ветви из
\[
|x-2|=4
\]
правильно возникают два линейных уравнения
\[
x-2=4
\quad\text{и}\quad
x-2=-4,
\]
то есть $x=6$ и $x=-2$. Первый неверный результат --- исключение $x=-2$ без основания. Затем вторая ветвь остаётся незавершённой.

\textbf{Почему это неверно.} Уравнение вида
\[
|A|=c,
\qquad c>0,
\]
имеет два случая:
\[
A=c
\quad\text{или}\quad
A=-c.
\]
Оба найденных значения нужно сохранять, если они удовлетворяют исходному уравнению. В данном задании внешнее раскрытие даёт две положительные правые части, $4$ и $2$, поэтому каждая из двух внутренних ветвей даёт по два корня.

\idea{} удобно раскрывать вложенный модуль слоями: сначала внешний модуль, затем отдельно решать каждое полученное уравнение с внутренним модулем и в конце объединить все корни.

\textbf{Исправление от последнего правильного шага.} Из
\[
|x-2|=4
\]
получаем
\[
x-2=4\Rightarrow x=6,
\]
\[
x-2=-4\Rightarrow x=-2.
\]
Оба значения сохраняем.

Из
\[
|x-2|=2
\]
получаем
\[
x-2=2\Rightarrow x=4,
\]
\[
x-2=-2\Rightarrow x=0.
\]

\textbf{Полное правильное решение.} Исходное уравнение:
\[
\bigl||x-2|-3\bigr|=1.
\]
Раскрываем внешний модуль:
\[
|x-2|-3=1
\quad\text{или}\quad
|x-2|-3=-1.
\]
Следовательно,
\[
|x-2|=4
\quad\text{или}\quad
|x-2|=2.
\]

Первая ветвь:
\[
|x-2|=4
\]
даёт
\[
x=6
\quad\text{или}\quad
x=-2.
\]

Вторая ветвь:
\[
|x-2|=2
\]
даёт
\[
x=4
\quad\text{или}\quad
x=0.
\]

Объединяем все значения:
\[
x\in\{-2,0,4,6\}.
\]

\textbf{Проверка.} Для $x=-2$ и $x=6$ внутренний модуль равен $4$, поэтому внешний модуль равен $|4-3|=1$. Для $x=0$ и $x=4$ внутренний модуль равен $2$, поэтому внешний модуль равен $|2-3|=1$. Все четыре значения подходят.

\result{} $x=-2$, $x=0$, $x=4$ или $x=6$.

\subsection*{Задача на закрепление}
Решите уравнение
\[
\bigl||x+1|-4\bigr|=2.
\]

\vfill
\begin{center}
\small Лёвин Артём Александрович эксклюзивно для Анны Трапезниковой
\end{center}

\end{document}
